% Self-contained source: latexmk -pdf writeup.tex
\documentclass[11pt,a4paper]{article}
\usepackage[T1]{fontenc}
\usepackage[utf8]{inputenc}
\usepackage{lmodern}
\usepackage[margin=27mm]{geometry}
\usepackage{amsmath,amssymb,amsthm,mathtools}
\usepackage{microtype,enumitem,booktabs,array}
\usepackage{xcolor,tikz}
\usepackage{aliascnt}
\usetikzlibrary{arrows.meta,positioning,calc,fit,backgrounds}
\usepackage[colorlinks=true,linkcolor=blue!55!black,citecolor=blue!55!black,
  urlcolor=blue!55!black,bookmarksnumbered=true]{hyperref}
\usepackage[capitalise,nameinlink,noabbrev]{cleveref}
\usepackage{caption}
\captionsetup{font=small,labelfont=bf}
\newtheorem{theorem}{Theorem}[section]
\newaliascnt{lemma}{theorem}
\newtheorem{lemma}[lemma]{Lemma}
\aliascntresetthe{lemma}
\newaliascnt{proposition}{theorem}
\newtheorem{proposition}[proposition]{Proposition}
\aliascntresetthe{proposition}
\newaliascnt{corollary}{theorem}
\newtheorem{corollary}[corollary]{Corollary}
\aliascntresetthe{corollary}
\newaliascnt{claim}{theorem}
\newtheorem{claim}[claim]{Claim}
\aliascntresetthe{claim}
\theoremstyle{definition}
\newaliascnt{definition}{theorem}
\newtheorem{definition}[definition]{Definition}
\aliascntresetthe{definition}
\theoremstyle{remark}
\newaliascnt{remark}{theorem}
\newtheorem{remark}[remark]{Remark}
\aliascntresetthe{remark}
\DeclareMathOperator{\Def}{def}
\DeclareMathOperator{\supp}{supp}
\newcommand{\Q}{\mathbb Q}
\newcommand{\R}{\mathbb R}
\newcommand{\N}{\mathbb N}
\newcommand{\HH}{\mathcal H}
\newcommand{\PP}{\mathcal P}
\newcommand{\budget}{\mathcal R}
\newcommand{\LOCAL}{\textsf{LOCAL}}
\newcommand{\ind}{\mathbf 1}
\newcommand{\avg}[2]{\left\langle #1\right\rangle_{#2}}
\definecolor{vertexblue}{HTML}{275D84}
\definecolor{witnessorange}{HTML}{B75E17}
\definecolor{pairpurple}{HTML}{794581}
\tikzset{
  vertex/.style={circle,draw=vertexblue,fill=blue!8,thick,inner sep=2pt,minimum size=6mm},
  witness/.style={circle,draw=witnessorange,fill=orange!15,thick,inner sep=2pt,minimum size=6mm},
  edge/.style={thick},
  newedge/.style={very thick,draw=witnessorange},
  flowbox/.style={draw=vertexblue,rounded corners=2pt,fill=blue!4,
    align=center,font=\small,text width=31mm,minimum height=15mm},
  flowarrow/.style={-{Stealth[length=2mm]},thick,draw=vertexblue}
}
\hypersetup{
  pdftitle={An exact factor-seven approximation for planar domination in constant locality},
  pdfauthor={AI-generated write-up}
}
\title{An exact factor-seven approximation for planar domination\\
  in constant locality}
\author{AI-generated write-up}
\date{23 October 2026}

\begin{document}
\maketitle

\begin{abstract}
We prove that minimum dominating set in finite simple planar graphs admits a
deterministic constant-round \LOCAL{} algorithm with approximation ratio
exactly \(7\). The constant is independent of the number of vertices and the
maximum degree; the algorithm needs neither of these parameters nor a planar
embedding. Together with the known lower bound excluding every ratio
\(7-\varepsilon\), this shows that the optimal constant is attained.
The proof combines a two-threshold reduction of hard common neighborhoods,
exact maximal packing with bounded constraint size, a selection threshold
that retains quantitative slack, and a planar face-incidence count.
The comparison is with the integral domination number. We give the complete
algorithm and all the combinatorial and numerical arguments needed for its
guarantee, including a finite construction of its constant-round packing
routine.
\end{abstract}

\section{The question and the result}
\label{sec:introduction}

A set \(D\) of vertices of a graph is \emph{dominating} if every vertex
belongs to \(D\) or has a neighbor in \(D\). The minimum cardinality of such
a set is the \emph{domination number} \(\gamma(G)\). We ask how well a
distributed algorithm can approximate \(\gamma(G)\) when every vertex must
decide whether to join the output after a constant number of communication
rounds.

Hilke, Lenzen, and Suomela proved that, for every fixed
\(\varepsilon>0\), no deterministic constant-round algorithm achieves
approximation ratio \(7-\varepsilon\) on planar graphs~\cite{HLS14}.
Recent work of Amiri and Siebertz~\cite{AS26} and
Picavet~\cite{Picavet26} gives ratio \(7+\varepsilon\) for every fixed
\(\varepsilon>0\). These bounds identify the infimum of the achievable
ratios as \(7\), but leave the endpoint question:
\[
  \textit{Is there one constant-round algorithm with ratio exactly \(7\)?}
\]
The answer is affirmative.

\begin{theorem}\label{thm:main}
There are a deterministic \LOCAL{} algorithm and an absolute constant \(r\)
such that, on every finite simple planar graph \(G\), the algorithm terminates
within \(r\) rounds and outputs a dominating set \(O\) with
\[
  |O|\leq 7\gamma(G).
\]
The round bound is independent of \(|V(G)|\), the maximum degree, the
identifier assignment, and the port numbering. The algorithm is not given
a planar embedding, \(|V(G)|\), or the maximum degree.
\end{theorem}

The guarantee has no additive error and no minimum-order restriction.
In particular, disconnected graphs, isolated vertices, and the empty graph
are included. The constants in the construction are deliberately generous;
no useful practical running time is claimed.

\subsection{Prior work and the techniques used here}

In the deterministic LOCAL model~\cite{Linial92,Suomela13}, the established
constant-round approximation ratios for planar minimum dominating set have
progressed from $130$~\cite{LPW13}, through $52$~\cite{Waw14} and
$11+\varepsilon$~\cite{HSV22,HKSS25}, to $8+\varepsilon$~\cite{Waw26} and,
most recently, $7+\varepsilon$~\cite{AS26,Picavet26}. Hilke, Lenzen, and
Suomela~\cite{HLS14} proved that every ratio $7-\varepsilon$ is impossible in
constant rounds. Amiri and Siebertz~\cite{AS26} also give a path-power
lower-bound construction, including a formulation with known network size
and identifiers drawn from $\{1,\ldots,n\}$. These bounds identify seven as
the infimum; attaining that infimum requires an additional argument. In
other local optimization problems, an endpoint approached by all
$\alpha+\varepsilon$ guarantees can remain unattainable~\cite{FHK11}.

The two recent $7+\varepsilon$ algorithms have different scope and techniques.
Amiri and Siebertz~\cite{AS26} use clustering, averaging, and LP
derandomization on proper minor-closed classes; their planar guarantee is
against the fractional domination optimum $\gamma_f(G)$. Picavet~\cite{Picavet26}
works under bounds on maximum average degree and on the density of
depth-one shallow minors, with a guarantee against the integral optimum
$\gamma(G)$. Maximum average degree means the largest average degree of
a nonempty subgraph. A depth-one shallow minor is obtained by contracting
disjoint connected subgraphs of radius at most one and then deleting
vertices and edges; for a nonempty minor, its density is its number of
edges divided by its number of vertices. The exact-seven result here also
compares to $\gamma(G)$.
It does not establish a factor seven against $\gamma_f(G)$.
Here $\gamma_f(G)$ is the minimum of $\sum_{v\in V}x_v$ over
nonnegative real weights satisfying $\sum_{v\in N[w]}x_v\geq1$ for
every vertex $w$, where $N[w]$ is the vertex $w$ together with its
neighbors; integral dominating sets are feasible zero--one
solutions, so $\gamma_f(G)\leq\gamma(G)$.

The present construction starts from Picavet's hard-neighborhood hypergraph
and sunflower framework~\cite{Picavet26,ER60}, which in turn builds on the
degree-reduction and LP approach of Heydt and coauthors~\cite{HSV22,HKSS25}.
We also use the threshold/layer-cake viewpoint in Picavet's rounding
analysis. The endpoint proof separates two structural thresholds, computes
an exact maximal fractional packing without a bound on variable frequency,
and charges missed hard neighborhoods to quantitative rounding slack and
planar Euler deficits. Every lemma used for these steps is proved here.

Density-based domination rounding predates these local algorithms.
The arboricity $a$ of a graph is the minimum number of forests whose
edge sets cover all of its edges.
Bansal and Umboh~\cite{BU17} give a $3a$ approximation for graphs of
arboricity at most a positive bound $a$,
and Dvo\v r\'ak~\cite{Dvorak19} improves the relevant rounding bound to
$2a+1$. Dory, Ghaffari, and Ilchi~\cite{DGI24} obtain
$(2a+1)(1+\varepsilon)$ in degree-dependent distributed time. Earlier
constant-round approaches use local approximate covering-LP solvers of
Kuhn, Moscibroda, and Wattenhofer~\cite{KMW06,KMW16} after reducing the
incidence parameters. Our packing lemma instead needs exact maximality,
not an optimal packing objective. Exact fractional-packing algorithms have
precedents in \AA strand and Suomela~\cite{AS10}; the bounded-row-arity,
finite-capacity formulation used here is proved directly.

The endpoint does not follow by substituting \(\varepsilon=0\) in a
\((7+\varepsilon)\)-approximation: the number of rounds of that family may
diverge as \(\varepsilon\) tends to zero. The proof below instead fixes all
parameters once and accounts for the apparent approximation losses using
nonnegative slack in the same solution.

\subsection{Proof outline}

First, we mark single vertices and pairs whose common neighbors cannot be
dominated cheaply from outside the marked set. A reduction with two distinct
thresholds produces a hypergraph \(\HH\) of rank at most two and bounded
maximum degree. Every original hard set contains a retained hyperedge.
After removing the common neighbors of retained hyperedges from the
request set, each vertex has only a bounded number of residual neighbors.

Next, we compute a maximal fractional packing with variables \(y_w\) for
residual requests and \(z_e\) for retained hyperedges. Its rows have a fixed
maximum size, even though a variable can occur in arbitrarily many rows.
We prove that fixed rational capacities and bounded row size allow an
exact constant-round algorithm for this packing.

Finally, we select vertices whose packing rows are sufficiently close to
tight, and add the residual vertices still undominated. The planar
edge bound yields a quantitative rounding inequality. Comparison with
an arbitrary integral dominating set produces a common nonnegative
budget. A face count and star contractions show that hyperedges missed by
the comparison set consume only a small fraction of that budget. The
remaining slack gives the exact factor \(7\).

\begin{figure}[t]
\centering
\begin{tikzpicture}[node distance=5mm]
  \node[flowbox] (hard) {Hard singletons\\and pairs};
  \node[flowbox,right=of hard] (packing) {Residual requests\\and packing \((y,z)\)};
  \node[flowbox,right=of packing] (threshold) {Near-tight rows\\select \(S\)};
  \node[flowbox,right=of threshold] (repair) {Repair uncovered\\requests with \(U\)};
  \draw[flowarrow] (hard)--(packing);
  \draw[flowarrow] (packing)--(threshold);
  \draw[flowarrow] (threshold)--(repair);
\end{tikzpicture}
\caption{The algorithm. All four stages use a fixed-radius neighborhood.
The marked hyperedges help dominate the removed requests; the repair set
\(U\) handles the residual requests not already dominated by \(S\).}
\label{fig:pipeline}
\end{figure}

\section{Model, notation, and elementary planar facts}
\label{sec:preliminaries}

We use the standard synchronous deterministic \LOCAL{} model
\cite{Linial92,Suomela13}. Vertices are processors with distinct positive
integer identifiers.
An edge permits a message in each direction in each round; messages and
local computations have no size or time bound. Incident edges may have
arbitrary local port numbers. The algorithm must be effective: every local
step is an ordinary terminating computation. Exact operations on rational
numbers and exhaustive searches over a finite neighborhood are permitted.
In a constant number of rounds a vertex can gather a fixed-radius
neighborhood, together with identifiers and the adjacency information
needed below. The graph is promised to be planar; the algorithm does not
test this promise.

All graphs are finite, undirected, and simple. Unless specified otherwise,
adjacency and neighborhoods refer to the input graph \(G=(V,E)\). Write
\[
  N(v)=\{w:vw\in E\},\qquad N[v]=N(v)\cup\{v\},\qquad
  N[A]=\bigcup_{a\in A}N[a].
\]
For a nonempty set \(A\subseteq V\), let
\[
  X(A)=\bigcap_{a\in A}N(a)
\]
be its common \emph{open} neighborhood. A set \(C\) dominates a request
set \(W\) if \(W\subseteq N[C]\). In particular, \(X(\{v\})=N(v)\).
For \(W\subseteq V\), put
\[
  e_G(W)=|E(G[W])|,
  \qquad
  \Def_G(W)=3|W|-e_G(W).
\]
We omit the subscript \(G\) when the ambient graph is clear. The
\emph{coarse deficit} \(\Def(W)\), with \(3|W|\) rather than
\(3|W|-6\), is convenient for small and disconnected induced subgraphs.
For a function \(y:V\to\R\), write \(y(W)=\sum_{w\in W}y_w\).
All sums over empty sets are zero.

\subsection{Planar drawings, edge bounds, and contractions}

A graph is planar if it has a drawing on the sphere in which vertices
are distinct points and edges are arcs meeting only at shared endpoints.
The following standard facts about planar drawings are included with
proofs; see also Diestel~\cite[Chapter~4]{Diestel25} for general background.
A planar drawing is used only in the analysis, not supplied to the algorithm.

For disjoint $A,B$, let $e_G(A,B)$ count edges with one endpoint in each
set. The graph consisting of these edges is bipartite, regardless of
any other edges of $G$.

Fix a drawing of a connected graph with at least one edge. A \emph{dart}
is an oriented edge $(u,v)$, and its reverse is $(v,u)$. The drawing gives
a cyclic order of the darts leaving each vertex. Let $\sigma$ take each
dart to its successor in this order, and put $\alpha(u,v)=(v,u)$.
The cycles of
\[
 \varphi=\sigma\circ\alpha
\]
are the facial boundary walks, traversed in one consistent direction.
Every dart belongs to exactly one such cycle, so the sum of boundary
lengths is $2e$. These walks may repeat vertices and may traverse both
orientations of a bridge; they need not be simple cycles.

For a connected drawing with $n$ vertices, $e$ edges, and $f$ faces,
\begin{equation}\label{geom:Euler}
 n-e+f=2.
\end{equation}
For completeness, an embedded tree has one face: successively deleting
leaves does not change the number of complementary regions and reduces
to a single vertex. Since a tree has $n-1$ edges, the formula holds.
Otherwise, delete an edge on a cycle. The graph stays connected, and the
two sides of the edge belong to distinct faces, since a simple cycle
containing it separates the sphere. Deletion merges precisely those two
faces, reducing both $e$ and $f$ by one. Repeating until a tree remains
proves \eqref{geom:Euler}.

\begin{lemma}[Planar edge bounds]\label{geom:edge-bounds}
A connected simple planar graph with $n\geq3$ vertices has at most $3n-6$
edges. If it is bipartite, it has at most $2n-4$ edges. Consequently every
finite simple planar graph, with no connectivity or minimum-order
assumption, satisfies
\[
 |E(G)|\leq3|V(G)|,\qquad
 |E(G)|\leq2|V(G)|\quad\text{if $G$ is bipartite}.
\]
In particular, $\Def(W)\geq0$ for every vertex set $W$, and
$e_G(A,B)\leq2(|A|+|B|)$ for disjoint $A,B$.
\end{lemma}
\begin{proof}
A facial boundary cannot have length one, because that would give a loop.
A boundary of length two has darts $(u,v),(v,u)$: successive darts form a
walk, and simplicity permits only this return in two steps. The equalities
\[
 \varphi(u,v)=(v,u),\qquad \varphi(v,u)=(u,v)
\]
say that the cyclic successor fixes $(v,u)$ at $v$ and $(u,v)$ at $u$.
A cyclic permutation fixes a member only if its fiber has size one.
Hence both endpoints have degree one, so their connected component is
$K_2$. Every facial boundary in a connected graph of order at least three
therefore has length at least three. Thus $3f\leq2e$, which together with
\eqref{geom:Euler} gives $e\leq3n-6$.

In a bipartite graph every closed walk has even length, since each step
changes sides. For a connected component of order at least three, the
preceding argument excludes length two, so every face has length at least
four. The inequality $4f\leq2e$ and \eqref{geom:Euler} give $e\leq2n-4$.

Components of order one and two have respectively zero and at most one
edge. They satisfy both relaxed bounds. Apply the connected bounds to
all other components and sum; the empty graph also satisfies the claims.
Induced subgraphs, and the bipartite subgraph between $A$ and $B$, inherit
planarity by deleting vertices and edges. This proves the final claims.
\end{proof}

\begin{lemma}[Outside domination of a common neighborhood]
\label{geom:outside-pair}
For three distinct vertices $a,b,d$ of a simple planar graph,
\[
 |N(a)\cap N(b)\cap N[d]|\leq3.
\]
\end{lemma}
\begin{proof}
The graph $K_{3,3}$ is not planar: it is connected and bipartite, with
six vertices and nine edges, whereas Lemma~\ref{geom:edge-bounds} allows
at most eight. If the displayed intersection had at least four members,
choose three distinct members $w_1,w_2,w_3$ different from $d$. Each is
adjacent to $a,b,d$. None equals $a$ or $b$, by absence of loops, and none
equals $d$ by choice. The two triples therefore form a $K_{3,3}$ subgraph.
Deleting all other vertices and edges would preserve its planar drawing,
a contradiction.
\end{proof}

\begin{lemma}[Simultaneous contraction of domination stars]
\label{geom:star-contraction}
Let $D$ dominate a planar graph $G$, and fix a map $f:V(G)\to D$ with
\[
 f(v)\in N[v]\quad(v\in V(G)),\qquad f(d)=d\quad(d\in D).
\]
For any $L\subseteq V(G)\setminus D$, put
\[
 q(v)=
 \begin{cases} f(v),&v\in L,\\ v,&v\notin L.\end{cases}
\]
The simple graph $F$ on $V(G)\setminus L$ with edge set
\[
 \bigl\{\{q(u),q(v)\}:uv\in E(G),\ q(u)\ne q(v)\bigr\}
\]
is planar. All vertices outside $L$ remain distinct. If $W\cap L$ is
empty, every edge of $G[W]$ remains in $F[W]$.
\end{lemma}
\begin{proof}
Such an assignment exists by fixing $D$ pointwise and choosing a dominating
neighbor in $D$ for every other vertex. If $v\in L$, then $v\ne f(v)$ and
the edge $vf(v)$ exists. These edges form vertex-disjoint stars centered
in $D$: every leaf has one assigned center, and no center is a leaf.

Contract all these star edges, retaining each center's name. Edge
contraction preserves a planar drawing. Indeed, choose a narrow disk
around an edge and its endpoints, containing no other vertex. The
incident arcs can be redrawn inside this disk to meet one vertex in their
existing boundary order, without crossings. Carry this out one edge at
a time. Delete resulting loops, and keep one copy of each group of
parallel edges. Neither deletion introduces a crossing. No two centers,
and no two vertices outside $L$, are identified. The resulting vertex
map is exactly $q$, and its simple edge set is the one stated.
For $u,v\in W$ the map fixes both distinct endpoints, so each edge $uv$
survives.
\end{proof}

\begin{lemma}[Realizing assigned incidences]\label{geom:assigned-incidences}
Keep $D,f$ as above. Let $A\subseteq V(G)\setminus D$, $B\subseteq D$,
and $P\subseteq A\times B$ be a set of distinct ordered pairs. Suppose
each $(a,d)\in P$ has a witness $w_{a,d}$ with
\[
 aw_{a,d}\in E(G),\qquad f(w_{a,d})=d,\qquad w_{a,d}\notin A.
\]
Then
\begin{equation}\label{geom:assigned-bipartite}
 |P|\leq2(|A|+|B|).
\end{equation}
If also $A\cup B\subseteq W$, all witnesses lie outside $W$, and
$ad\notin E(G)$ for every $(a,d)\in P$, then
\begin{equation}\label{geom:assigned-new}
 |P|\leq\Def(W).
\end{equation}
\end{lemma}
\begin{proof}
Let $L$ consist of the chosen witnesses outside $D$, and apply
Lemma~\ref{geom:star-contraction}. No vertex of $A\cup B$ is contracted.
Every witness maps to its assigned vertex $d$. This includes a witness
already in $D$: it is then equal to $f(w)=d$ and is left fixed.
The edge $aw_{a,d}$ consequently becomes $ad$, with distinct endpoints
because $A\cap D$ is empty.

Different ordered pairs in $P$ give different unordered edges: disjointness
of $A,B$ determines the orientation from $A$ to $B$. Restricting to those
edges gives a simple planar bipartite graph on $A\cup B$. Its edge bound
proves \eqref{geom:assigned-bipartite}.

Under the additional assumptions, the contractions fix $W$, so $F[W]$
contains every original edge of $G[W]$. The $|P|$ edges obtained above are
all distinct and were absent from $G[W]$. Hence
\[
 e_G(W)+|P|\leq |E(F[W])|\leq3|W|,
\]
proving \eqref{geom:assigned-new}.
\end{proof}

\begin{figure}[tbp]
\centering
\begin{tikzpicture}[scale=1.05]
 \coordinate (a) at (0,1);
 \coordinate (b) at (0,-1);
 \coordinate (d) at (2.4,0);
 \coordinate (w1) at (1.15,0.85);
 \coordinate (w2) at (1.15,-0.85);
 \draw[edge] (a)--(b) (a)--(w1) (b)--(w2);
 \draw[newedge] (w1)--(d) (w2)--(d);
 \node[vertex] at (a) {\(a\)};
 \node[vertex] at (b) {\(b\)};
 \node[vertex] at (d) {\(d\)};
 \node[witness] at (w1) {\(w_1\)};
 \node[witness] at (w2) {\(w_2\)};
 \node[font=\small,align=center] at (1.2,-1.65)
   {Before: \(f(w_1)=f(w_2)=d\)};
 \draw[flowarrow] (3.1,0)--(4.4,0)
   node[midway,above,font=\small] {contract};
 \coordinate (aa) at (5.1,1);
 \coordinate (bb) at (5.1,-1);
 \coordinate (dd) at (7.5,0);
 \draw[edge] (aa)--(bb);
 \draw[newedge] (aa)--(dd) (bb)--(dd);
 \node[vertex] at (aa) {\(a\)};
 \node[vertex] at (bb) {\(b\)};
 \node[vertex] at (dd) {\(d\)};
 \node[font=\small,align=center] at (6.3,-1.65)
   {After: old edge \(ab\) is retained};
\end{tikzpicture}
\caption{Contraction of a domination star. Blue vertices are retained,
and the orange witness vertices are contracted to their assigned
dominator \(d\). If the blue vertices lie in a protected set \(W\)
and the witnesses lie outside it, every original edge of \(G[W]\)
survives. Distinct new edges inside \(W\) can then be charged to
\(\Def(W)=3|W|-e_G(W)\). The picture illustrates the operation, without
assuming that this small example itself has a hard neighborhood.}
\label{fig:contraction}
\end{figure}

\begin{lemma}[Pair counts from endpoint degrees]\label{geom:pair-degree}
Let $\mathcal P$ be a simple graph of pairs of maximum degree at most
$T$. For any vertex set $Z$, at most $T|Z|/2$ of its edges have both
endpoints in $Z$. If all its endpoints lie in $A\subseteq V(G)$, put
\[
 c(w)=|\{\{a,b\}\in E(\mathcal P):a,b\in N_G(w)\}|.
\]
For planar $G$,
\begin{align}
 \sum_{w\in A}c(w)&\leq3T|A|,\label{geom:pairs-inside}\\
 \sum_{w\in B}c(w)&\leq T(|A|+|B|)
       \quad\text{if }A\cap B=\varnothing.\label{geom:pairs-bipartite}
\end{align}
The edges of $\mathcal P$ need not be edges of $G$.
\end{lemma}
\begin{proof}
The sum of degrees in $\mathcal P[Z]$ counts every internal pair twice
and is at most $T|Z|$. Apply this to $Z=N_G(w)\cap A$. For $w\in A$,
summing gives
\[
 \sum_{w\in A}c(w)
 \leq\frac T2\sum_{w\in A}|N_G(w)\cap A|
 =T e_G(A)\leq3T|A|.
\]
For disjoint $A,B$, it gives
\[
 \sum_{w\in B}c(w)
 \leq\frac T2 e_G(A,B)\leq T(|A|+|B|).
\]
Both last inequalities use Lemma~\ref{geom:edge-bounds}.
\end{proof}

\section{Exact maximal packing with bounded row size}
\label{sec:packing}

Maximal fractional packings are a standard tool for local covering
algorithms; see, for example, \AA strand and Suomela~\cite{AS10}.
We need an exact packing result whose round bound depends on constraint
size and on a fixed finite set of capacities, but not on the number of
constraints containing a variable.  We prove that result here.

Let $I$ be a finite set of variables and $\mathcal A$ a finite set of row
indices.  A row $r\in\mathcal A$ has a support $E_r\subseteq I$ and a
capacity $c_r>0$.  Different rows may have the same support and capacity.
There are no repeated occurrences of a variable within a row.  Given a
common cap $M>0$, a vector $x\in\mathbb R^I$ is \emph{feasible} if
\begin{equation}\label{pack:feasible}
 0\leq x_i\leq M\quad(i\in I),
 \qquad \sum_{i\in E_r}x_i\leq c_r\quad(r\in\mathcal A).
\end{equation}
It is \emph{coordinatewise maximal} if there is no other feasible vector
$x'$ with $x'_i\geq x_i$ for every $i$ and a strict inequality for at least
one $i$.  In the incidence network there is one node for each variable and
one node for each row, with an edge $ir$ exactly when $i\in E_r$.
Each row node knows its capacity.  Empty rows are harmless and may be
discarded.  We allow isolated variable nodes.

\begin{lemma}[Exact bounded-arity packing]\label{pack:theorem}
Fix an integer $K\geq0$, a finite set $C\subseteq\mathbb Q_{>0}$, and
$M\in\mathbb Q_{>0}$.  There is an effectively constructible finite set
\[
 \mathcal P(K,C,M)\subseteq[0,M]\cap\mathbb Q,
 \qquad 0,M\in\mathcal P(K,C,M),
\]
with the following property.  On every finite incidence network satisfying
$|E_r|\leq K$ and $c_r\in C$, a deterministic algorithm computes a
coordinatewise maximal feasible vector using at most
\[
 2\bigl(|\mathcal P(K,C,M)|-1\bigr)
\]
communication rounds.  Every output coordinate
belongs to $\mathcal P(K,C,M)$.  The algorithm needs neither the size of
the instance nor a bound on the number of rows containing a variable.
It uses exact rational arithmetic and does not require identifiers or
symmetry breaking on the incidence network.
\end{lemma}

\subsection{A certificate and a finite progressive-filling process}

We call a feasible vector $x$ a \emph{water-filling certificate} if, for
each $i\in I$, either $x_i=M$, or there is a row $r$ such that
\begin{equation}\label{pack:certificate}
 i\in E_r,\qquad
 \sum_{j\in E_r}x_j=c_r,\qquad
 x_j\leq x_i\quad\hbox{for every }j\in E_r.
\end{equation}
Thus a noncapped coordinate has a tight row whose other coordinates are
no larger.  This last condition is the additional information used in the
palette argument.

\begin{lemma}\label{pack:existence}
Every finite system with positive rational capacities and positive
rational common cap has a rational water-filling certificate.
Every water-filling certificate is coordinatewise maximal.
\end{lemma}
\begin{proof}
Here is a finite, global process used to prove existence.  Its global
minimum operations will not be part of the distributed algorithm.
Initially all coordinates are active, their values are zero, and the
current level is $a=0$.  Active coordinates always have the common value
$a$; every already frozen coordinate has its final value, at most $a$.

Suppose the active set $A$ is nonempty.  For each row put
\[
 m_r=|E_r\cap A|,
 \qquad F_r=\sum_{i\in E_r\setminus A}x_i.
\]
For rows with $m_r>0$ define
$b_r=(c_r-F_r)/m_r$, and let
\begin{equation}\label{pack:next-event}
 b=\min\bigl(\{M\}\cup\{b_r:r\in\mathcal A,\ m_r>0\}\bigr).
\end{equation}
The set in this minimum is finite and nonempty, and all its entries are
rational.  Current feasibility gives $b_r\geq a$ for every such row,
and hence $a\leq b\leq M$.  Raise every active coordinate to $b$ and
leave frozen coordinates unchanged.  A row with $m_r>0$ now has load
$F_r+m_rb\leq c_r$; a row with $m_r=0$ has not changed.  Thus the new
vector is feasible.

Freeze every active coordinate lying in any row that is now tight;
also freeze all active coordinates if $b=M$.  All simultaneous tight
rows are processed in this step.  At least one active coordinate freezes:
either $M$ attains the minimum and all of them freeze, or some row with
$m_r>0$ attains the minimum and all of its active coordinates freeze.
For a coordinate frozen by a tight row, all coordinates of that row are
now frozen, and their values are at most $b$.  Consequently that row
stays tight forever and supplies~\eqref{pack:certificate} for the
coordinate just frozen.  Previously supplied certificates are preserved.
The level is now $b$, and the process repeats on the remaining active set.

This proves termination after at most $|I|$ steps.  More explicitly, when
any active coordinates remain, no tight row contains an active
coordinate, since the preceding step froze all such coordinates.  Also
the current level is below $M$.  Hence the next level in
\eqref{pack:next-event} is strictly larger than the current level.
The initial positivity assumptions give the same conclusion for the
first step.  Every operation uses rational arithmetic, so the final
vector is rational.  If $I$ is empty there is no step and the conclusion
is vacuous.  An isolated variable is eventually frozen at $M$.
The final vector has the asserted certificate.

Finally, let $x$ be any vector satisfying the certificate, and suppose
$x'\geq x$ is feasible.  If $x'_i>x_i$, then $x_i\ne M$.  A tight row
witnessing~\eqref{pack:certificate} for $i$ has strictly larger load at
$x'$ than at $x$, contradicting feasibility.  Thus no coordinate can
increase, as required.
\end{proof}

\subsection{An explicit finite palette}

We now define $\mathcal P(K,C,M)$ recursively.  This is a finite
computation on rational numbers; it does not examine the incidence
network.  For $K=0$, or for $C=\varnothing$, set
\begin{equation}\label{pack:base-palette}
 \mathcal P(K,C,M)=\{0,M\}.
\end{equation}
For $K\geq1$ and $C\ne\varnothing$, form the primary set
\begin{equation}\label{pack:primary}
 B=\{0,M\}\cup
 \left\{\frac{c}{h}:c\in C,\ 1\leq h\leq K,
                    \ 0<\frac{c}{h}<M\right\}.
\end{equation}
Write its distinct elements in increasing order as
$0=b_0<b_1<\cdots<b_s=M$.  Initialize $F=\{0\}$.
Process the consecutive intervals $(a,b)=(b_j,b_{j+1})$ from left to
right.  At the beginning of an interval define the finite set
\begin{equation}\label{pack:residual-capacities}
 C_{a,F}=\mathbb Q_{>0}\cap\left\{
 c-\sum_{\nu=1}^{d}p_\nu-ma:
 \begin{array}{l}
 c\in C,\quad d,m\geq1,\quad d+m\leq K,\\
 p_1,\ldots,p_d\in F
 \end{array}\right\}.
\end{equation}
The notation in this display means that only positive values of the
displayed expression are retained.  Repetition among the $p_\nu$ is
allowed.  Compute
\[
 P'=\mathcal P(K-1,C_{a,F},b-a)
\]
and replace $F$ by
\begin{equation}\label{pack:palette-update}
 F\ \longleftarrow\ F\cup\{a+t:t\in P'\}\cup\{b\}.
\end{equation}
After the last interval, return $F$ as $\mathcal P(K,C,M)$.

There are finitely many primary intervals and finitely many choices in
\eqref{pack:residual-capacities}.  Every recursive call decreases $K$;
its new common cap $b-a$ is positive and its capacity set is finite and
positive.  The procedure therefore terminates.  All comparisons,
sorting, and arithmetic are exact operations on rational numbers, so
the palette is effectively constructible.  By induction, the recursive
palette $P'$ lies in $[0,b-a]$, and its translate lies in $[a,b]$.
At the start of an interval, $F\subseteq[0,a]$; after the update,
$F\subseteq[0,b]$.  It follows in particular that
$\mathcal P(K,C,M)\subseteq[0,M]\cap\mathbb Q$, and that the final
palette contains both $0$ and $M$.  These facts also cover an empty
residual capacity set.  No quantity in this construction depends on
$|I|$, $|\mathcal A|$, or variable frequency.

\begin{lemma}[Palette coverage]\label{pack:coverage}
If $|E_r|\leq K$, $c_r\in C$, and $x$ is a water-filling certificate
with common cap $M$, then
\[
 x_i\in\mathcal P(K,C,M)\quad\hbox{for every }i\in I.
\]
\end{lemma}
\begin{proof}
We prove the assertion by induction on $K$, for every choice of finite
positive rational $C$ and positive rational $M$, and every finite system.
If $K=0$, no row can contain a variable.  If $C=\varnothing$, there are
no rows at all.  In either case the certificate forces $x_i=M$ for every
$i$, proving the assertion for the base cases.

Assume $K\geq1$, $C\ne\varnothing$, and the assertion for $K-1$.
In the primary-interval sweep maintain the invariant
\begin{equation}\label{pack:sweep-invariant}
 x_i\leq a\quad\Longrightarrow\quad x_i\in F.
\end{equation}
It holds initially at $a=0$, since feasibility makes every such value
zero.  Consider one consecutive primary interval $(a,b)$.  By the
definition of the primary set, for every $c\in C$ and every
$1\leq h\leq K$ we have
\begin{equation}\label{pack:gap}
 \frac{c}{h}\leq a\quad\hbox{or}\quad\frac{c}{h}\geq b.
\end{equation}
Indeed, a ratio strictly between $a$ and $b$ would be an omitted primary
point.  Ratios at the boundaries are allowed in this statement; ratios
at least $M$ satisfy the second alternative.

Let $I_a=\{i\in I:x_i>a\}$.  For each row split its support into
\[
 O_r=\{i\in E_r:x_i\leq a\},
 \qquad L_r=\{i\in E_r:x_i>a\}.
\]
Keep only rows for which $O_r$ and $L_r$ are both nonempty.  For every
kept row define a residual row on $I_a$ with support $L_r$ and capacity
\begin{equation}\label{pack:residual-row}
 c'_r=c_r-\sum_{i\in O_r}x_i-|L_r|a.
\end{equation}
Its arity is at most $K-1$, because at least one old coordinate was
deleted.  Its capacity is positive: feasibility and $L_r\ne\varnothing$
give
\[
 c'_r\ \geq\ \sum_{i\in L_r}(x_i-a)\ >\ 0.
\]
By the sweep invariant every old value belongs to $F$.  Taking
$d=|O_r|$, $m=|L_r|$ in~\eqref{pack:residual-capacities} therefore
shows that $c'_r\in C_{a,F}$.  In particular the residual capacities
are rational, even if we initially allowed a real-valued certificate.

Give this residual system the common cap $b-a$, and set
\begin{equation}\label{pack:residual-vector}
 z_i=\min\{x_i,b\}-a\quad(i\in I_a).
\end{equation}
We claim that $z$ is a water-filling certificate for the residual
system.  First, $0<z_i\leq b-a$, and each retained row satisfies
\begin{align*}
 \sum_{i\in L_r}z_i
 &=\sum_{i\in L_r}\min\{x_i,b\}-|L_r|a\\
 &\leq c_r-\sum_{i\in O_r}x_i-|L_r|a
 =c'_r.
\end{align*}
Thus $z$ is feasible.

If $x_i\geq b$, then $z_i=b-a$, so the common-cap alternative of
the certificate holds.  It remains to consider $a<x_i<b$.  Since
$b\leq M$, this coordinate is not capped in the original system.
Choose a witness row $r$ from~\eqref{pack:certificate}, and put
$h=|E_r|$.  It contains $i$, so $1\leq h\leq K$, and
\[
 \frac{c_r}{h}
 =\frac{1}{h}\sum_{j\in E_r}x_j
 \leq x_i<b.
\]
Equation~\eqref{pack:gap} gives $c_r/h\leq a$.
If $O_r$ were empty, all $h$ coordinates of this row would be strictly
greater than $a$, and its tightness would yield
$c_r>ha\geq c_r$, a contradiction.  Hence $O_r\ne\varnothing$;
also $i\in L_r$, so this row was retained.
Every original coordinate in the witness row is at most $x_i<b$.
No one of them is altered by truncation at $b$, and therefore
\[
 \sum_{j\in L_r}z_j
 =\sum_{j\in L_r}x_j-|L_r|a=c'_r.
\]
Furthermore, $z_j=x_j-a\leq x_i-a=z_i$ for every $j\in L_r$.
This is precisely the required residual certificate.

The induction hypothesis now applies to the entire residual system,
including the case of no residual rows or no residual variables.  Thus
$z_i\in\mathcal P(K-1,C_{a,F},b-a)$ for each $i\in I_a$.
If $a<x_i<b$, then $x_i=a+z_i$ and the update
\eqref{pack:palette-update} lists $x_i$.  Values $x_i\leq a$ remain in
$F$, and the boundary value $x_i=b$ is added explicitly.  We have proved
the sweep invariant with $b$ in place of $a$.  After the final interval
we have $b=M$, and feasibility gives $x_i\leq M$ for every coordinate.
This completes the induction.
\end{proof}

Two features of the interval argument are useful to make explicit.
A row all of whose coordinates have value at most $a$ has no residual
variable and cannot supply a new interior event.  A row with no such
coordinate cannot witness an event in $(a,b)$: the average-capacity
argument in the proof forces every witness row for such an event to
contain an old coordinate.  Thus only rows with both types are needed
in the residual system.  Also, the artificial cap $b-a$ only truncates
coordinates whose original final values are at least $b$.  It does not
assert that those original coordinates freeze at $b$.  A value equal to
$b$ is covered separately, and later intervals handle values above $b$.
These observations apply equally at the first interval, at the last
interval, and when some primary points coincide before duplicates are
removed.

\subsection{The fixed scan and its locality}

Sort the complete palette as
\[
 \mathcal P(K,C,M)=\{0=t_0<t_1<\cdots<t_L=M\}.
\]
The incidence-network algorithm begins with value zero at every variable
and with all variables unfrozen.  For each $h=1,\ldots,L$, it performs
the following two communication phases.
\begin{enumerate}
\item Every unfrozen variable changes its value to $t_h$; a frozen
variable keeps its current value.  Every variable sends its value to
each incident row.
\item Every row computes its exact load and reports to every incident
variable whether the row is tight.  A variable freezes if it receives
a tight report or if $t_h=M$.
\end{enumerate}
Variables already frozen remain frozen.  A variable with no incident
row receives no report and freezes at the common cap.  The scan uses a
fixed finite schedule independent of the instance; it does not wait for
a globally detected next event.

\begin{proof}[Proof of Lemma~\ref{pack:theorem}]
Let $x^*$ be a certificate supplied by Lemma~\ref{pack:existence}.
Every coordinate of $x^*$ is positive: if $x^*_i=0$, then it is below
$M$, and a witness row would have all its nonnegative coordinates at
most zero, contradicting its positive capacity.
Lemma~\ref{pack:coverage} gives $x^*_i\in\mathcal P(K,C,M)$.
We prove by induction on the scan levels that after level $t_h$ the
current value at $i$ is
\begin{equation}\label{pack:scan-invariant}
 \min\{x^*_i,t_h\},
 \qquad\hbox{and $i$ is frozen exactly when }x^*_i\leq t_h.
\end{equation}
At $t_0=0$ this follows from positivity and initialization.
Suppose it holds after level $a=t_{h-1}$ and write $t=t_h$.
There is no coordinate value of $x^*$ strictly between $a$ and $t$,
because both are consecutive points of a palette containing all such
values.  A coordinate frozen at the previous level has its final value
$x^*_i\leq a$.  An unfrozen coordinate has $x^*_i>a$, hence
$x^*_i\geq t$, and is raised to $t$.  Thus after the first phase
the current vector is exactly $\min\{x^*,t\}$ coordinate by coordinate.
This vector is feasible, since it is nonnegative and coordinatewise no
larger than the feasible vector $x^*$.

If $x^*_i>t$, then $t<M$.  In every row containing $i$, the current
load is strictly smaller than the final load: this coordinate has been
strictly decreased, and no other coordinate has increased.  The final
load is at most capacity, so no incident row is tight now.  Thus $i$
does not freeze prematurely.
If $x^*_i=t$, then $i$ freezes at the common cap when $t=M$.
Otherwise, its witness row in~\eqref{pack:certificate} has every final
coordinate at most $t$.  The current values on that row already equal
their final values, and the row is tight, so $i$ freezes in this step.
Earlier-frozen coordinates stay frozen.  This proves the full invariant.

At the last level $M$, the output is exactly $x^*$.  It is rational,
feasible, and coordinatewise maximal by Lemma~\ref{pack:existence}.
The number of phases is $2L=2(|\mathcal P|-1)$.  Each incidence edge
carries one message in each direction per level.  A variable may have
arbitrarily many incident rows, but all its messages are sent
simultaneously.  The LOCAL model permits an arbitrarily large finite
message on every edge and unbounded local computation, so this does not
alter the round count.  No identifiers, orientation, or ordering of
incident edges enter the algorithm.  This proves the lemma.
\end{proof}

\subsection{Effective arithmetic}

The construction above gives a finite algorithm, rather than only the
existence of a suitable finite set of numbers.  The recursion for the
palette terminates by decreasing the nonnegative integer $K$, and its
inner loops enumerate finite sets and tuples of rational numbers.  It
can therefore be executed as a preliminary computation depending only
on the fixed parameters, or its finite output can be included in the
program.  Neither alternative requires information about the network.

To implement all scan operations using integers, take any positive
integer $D$ divisible by the denominators of every member of
$\mathcal P(K,C,M)\cup C\cup\{M\}$.  Such a $D$ is obtained, for
example, by multiplying the positive denominators in reduced fraction
representations.  Multiply every value, capacity, and scan level by
$D$.  Every variable then stores a nonnegative integer.  Row loads are
finite sums of integers, and tightness is an integer equality.  This
also shows that exact comparisons cause no issue of real-number
computation.  The size of the palette and of $D$ may be enormous; only
their finiteness and independence from the network are required.

\begin{remark}
The constant in \cref{pack:theorem} depends on row size and on a fixed
finite set of rational capacities; it need not be efficient in these
parameters. This is consistent with the \(\Omega(\Delta)\) locality
barrier for maximal fractional matching when the running time is
independent of the network size~\cite{GHS17}: its constraints
are indexed by vertices, its variables by edges, and its row size can
be the maximum degree \(\Delta\). The lemma makes no half-integrality
claim. Fine rational denominators can be necessary in local maximal
fractional matching~\cite{Dahal24}.
\end{remark}


\section{The hard-neighborhood reduction and the algorithm}
\label{sec:algorithm}

Fix the following constants for the rest of the paper:
\begin{equation}\label{eq:constants}
  \theta=\frac17,\qquad T=10^6,\qquad Q=10^{12}.
\end{equation}
In particular, \(1\leq T\leq Q\).

\subsection{Hard sets and the retained hypergraph}

A hypergraph here is simply a family of nonempty vertex sets, called
hyperedges. Its rank is the maximum hyperedge size, and a vertex's
hypergraph degree is the number of hyperedges containing it.


\begin{definition}
A nonempty set \(A\subseteq V\) is \emph{\(q\)-hard} if no set
\(C\subseteq V\setminus A\) with \(|C|\leq q\) dominates \(X(A)\).
\end{definition}

Let \(\HH_0\) consist of the \(Q\)-hard singletons and the \(Q\)-hard
two-element sets. For each vertex \(v\), put
\[
  P(v)=\{u\ne v:\{v,u\}\in\HH_0\}.
\]
Define the singleton centers and the retained hypergraph by
\begin{align}
  C_*&=\{v:\{v\}\in\HH_0\text{ or }|P(v)|>T\},\label{eq:centers}\\
  \HH&=\{\{v\}:v\in C_*\}
   \ \cup\
   \{e\in\HH_0:|e|=2,\ e\cap C_*=\varnothing\}.\label{eq:hypergraph}
\end{align}
Thus we add a singleton at every vertex incident with more than \(T\)
original hard pairs, and remove every pair that contains a singleton
center. The pairs sharing a center have disjoint one-vertex remainders,
the rank-two instance of a sunflower~\cite{ER60}. Its elementary
reduction properties are proved next.

Set
\begin{equation}\label{eq:residual}
  X=\bigcup_{e\in\HH}X(e),\qquad R=V\setminus X.
\end{equation}
The vertices of \(R\) are the \emph{residual requests}.

\begin{lemma}\label{lem:hard-properties}
The construction satisfies the following assertions.
\begin{enumerate}[label=(\roman*),itemsep=2pt]
  \item Every retained hyperedge has one or two vertices, every hard set
  has a nonempty common neighborhood, and the endpoints of every retained
  pair are at distance at most two.
  \item Every edge of \(\HH_0\) contains an edge of \(\HH\).
  Every retained singleton is \(T\)-hard, and every retained pair is
  \(Q\)-hard.
  \item Every vertex belongs to at most \(T\) hyperedges of \(\HH\).
  \item If \(1\leq |A|\leq2\) and \(X(A)\cap R\ne\varnothing\),
  then \(A\) is not \(Q\)-hard.
  \item For every \(v\in V\),
  \[
    |N(v)\cap R|\leq B_0,\qquad B_0:=Q(3Q+2).
  \]
\end{enumerate}
\end{lemma}
\begin{proof}
If \(X(A)=\varnothing\), the empty set dominates it, so \(A\) is not
hard. This proves the nonemptiness assertion and the distance bound for
pairs. A pair of \(\HH_0\) is either retained or contains a center of
\(C_*\); an original singleton is retained. Hence every original edge
contains a retained edge.

An original hard singleton is \(T\)-hard because \(T\leq Q\).
Suppose that \(v\) is a new center with \(|P(v)|>T\), and that
\(C\subseteq V\setminus\{v\}\), \(|C|\leq T\), dominates \(N(v)\).
For every \(u\in P(v)\), the set \(C\) dominates
\(X(\{v,u\})\subseteq N(v)\). If \(u\notin C\), this is a cover of that
common neighborhood from outside \(\{v,u\}\), contradicting its
\(Q\)-hardness. Thus \(P(v)\subseteq C\), which is impossible.
Retained pairs are original hard pairs.

A center belongs only to its own singleton in \(\HH\). A noncenter
belongs only to original pairs, of which it has at most \(T\).
For (iv), if \(A\) were \(Q\)-hard, it would contain a retained
hyperedge \(e\), and
\(X(A)\subseteq X(e)\subseteq X\), contrary to
\(X(A)\cap R\ne\varnothing\).

For (v), write \(m=|N(v)\cap R|\), and suppose \(m>Q(3Q+2)\).
By (iv) applied to \(\{v\}\), some set
\(C\subseteq V\setminus\{v\}\), \(|C|\leq Q\), dominates \(N(v)\).
For \(a\in C\), at most one vertex of \(N(v)\cap R\) is dominated by
equality with \(a\); all other vertices covered by \(a\) are in
\(N(v)\cap N(a)\cap R\). If these intersections all had at most \(3Q\)
vertices, their total cover would have size at most
\[
  |C|(3Q+1)\leq Q(3Q+1)<m.
 \]
Consequently some \(a\in C\) satisfies
\(|N(v)\cap N(a)\cap R|>3Q\). In particular \(a\ne v\), and
(iv) applied to \(\{v,a\}\) gives an outside cover of \(X(\{v,a\})\)
of size at most \(Q\). By the planar three-neighbor bound
(\cref{geom:outside-pair}), one vertex outside this pair dominates at most
three of its common neighbors. Such a cover cannot dominate more than
\(3Q\) common neighbors, a contradiction.
\end{proof}

\subsection{The mixed packing and the output}

For each \(w\in R\) introduce a variable \(y_w\), and for each
\(e\in\HH\) a variable \(z_e\). Compute a coordinatewise maximal feasible
vector for the following packing constraints:
\begin{align}
  \ell_v
  :=\sum_{w\in R\cap N[v]}y_w+
     \sum_{\substack{e\in\HH\\v\in e}}z_e
  &\leq 1 &&(v\in V),\label{eq:packing}\\
  0\leq y_w&\leq\theta &&(w\in R),\nonumber\\
  0\leq z_e&\leq1 &&(e\in\HH).\nonumber
\end{align}
The upper bounds \(z_e\leq1\) are redundant, because every hyperedge is
nonempty and occurs in the unit row of each endpoint. We state them
explicitly to use a common coordinate cap \(M=1\). The bounds
\(y_w\leq\theta\) are singleton rows of capacity \(\theta\).
By \cref{lem:hard-properties}, every unit row has at most
\[
  K:=B_0+1+T=Q(3Q+2)+1+T
\]
variables. All row capacities belong to \(\{1,1/7\}\).
Thus \cref{pack:theorem} supplies an exact maximal packing, with rational
coordinates, using a fixed finite list of levels.
Extend \(y\) by setting \(y_v=0\) for \(v\notin R\).

Define
\begin{equation}\label{eq:output}
  S=\left\{v:\ell_v\geq
     1-\frac{\theta-y_v}{2}\right\},
  \qquad
  U=R\setminus N[S],
  \qquad
  O=S\cup U.
\end{equation}
The threshold depends on \(y_v\). Selecting only the exactly tight rows
would lose a slack term that will be needed later.

\begin{lemma}\label{lem:output-basic}
The set \(O\) is dominating. Moreover, \(S\cap U=\varnothing\),
every \(e\in\HH\) meets \(S\), and
\begin{equation}\label{eq:ucap}
  y_w=\theta\quad\text{for every }w\in U.
\end{equation}
\end{lemma}
\begin{proof}
A tight unit row belongs to \(S\), since \(y_v\leq\theta\).
Consider a hyperedge \(e\). If \(z_e=1\), every endpoint has a tight
unit row. If \(z_e<1\) and no endpoint row is tight, then all rows
containing \(z_e\) have positive slack. There are finitely many such
rows, so \(z_e\) could be increased by some positive amount, contrary
to coordinatewise maximality. Thus \(e\cap S\ne\varnothing\).

Now let \(w\in U\). None of the unit rows containing \(y_w\) is tight:
such a row is indexed by a vertex of \(N[w]\), which would belong to
\(S\) and dominate \(w\). If \(y_w<\theta\), its personal cap and all
its incident rows therefore have positive slack, again allowing a
positive increase. This proves \eqref{eq:ucap}.

The definition of \(U\) gives \(S\cap U=\varnothing\) and ensures that
\(O\) dominates \(R\). For every \(w\in X\), there is an
\(e\in\HH\) with \(w\in X(e)\); any vertex in \(e\cap S\) is adjacent
to \(w\). Hence \(O\) also dominates \(X\), and \(R\cup X=V\).
\end{proof}

\subsection{An effective constant-round implementation}
\label{sec:implementation}

\paragraph{The preprocessing takes at most six rounds.}
We give an explicit schedule, including the initial exchange of identifiers.
In the first round, every vertex sends its identifier to all its neighbors.
In the second, it sends its complete list of neighbor identifiers.
After these rounds a vertex $v$ knows $N(v)$ and $N(w)$ for every
$w\in N(v)$.  In particular it knows every vertex at distance at most
two from $v$, and every edge incident with a neighbor of $v$.

This information determines all original hard sets containing $v$.
For a singleton the request set is $N(v)$.  For a candidate pair
$\{v,u\}$, it is
\[
 X(\{v,u\})=\{w\in N(v):u\in N(w)\},
\]
which $v$ can compute from the received lists.  Only vertices
$u\ne v$ at distance at most two need be considered: a pair with no
common neighbor is not hard, and every pair with a common neighbor
has this distance bound.  For either kind of tested set $A$ containing
$v$, put $W=X(A)\subseteq N(v)$.  Any vertex that dominates at least
one request of $W$ lies in
\[
 N[W]=\bigcup_{w\in W}N[w]\subseteq B_2(v),
\]
where $B_2(v)$ is the closed distance-two ball.  In a proposed cover,
vertices outside $N[W]$ cover no request and can be deleted.  Thus $v$
can decide $Q$-hardness by enumerating all subsets of
$N[W]\setminus A$ of size at most $Q$ and testing whether they cover
$W$.  The received neighbor lists give exactly the incidences needed
for every such test.  This also handles $W=\varnothing$, for which
the empty cover succeeds.  Consequently, by the end of round two,
$v$ knows its complete hard-partner set $P(v)$ and whether it belongs
to the singleton-center set $C_*$.

In rounds three and four, vertices forward their center bits, tagged
with their identifiers, to distance two.  Every $v$ then knows the
center bits of all $u\in P(v)$, so it can determine its complete list
of incident retained hyperedges: its singleton if $v\in C_*$, and
each original hard pair whose two endpoints lie outside $C_*$.  All
endpoints obtain consistent descriptions because both hardness and
the two-center test are symmetric.

In round five, every vertex sends its list of incident retained
hyperedges to all its graph neighbors.  A receiving vertex $w$ is in
$X$ exactly when one of the received hyperedges $e$ satisfies
$e\subseteq N(w)$.  The forward implication follows because every
endpoint of such an $e$ is a neighbor of $w$ and sends its incident
edge list; the reverse implication is the definition of $X(e)$.
Thus $w$ determines whether it belongs to $R=V\setminus X$.
In round six, vertices exchange their residual-membership bits.
Every vertex $v$ now knows $R\cap N[v]$, all retained hyperedges
incident with $v$, and hence its complete mixed packing row.

The length-one or length-two routes used for pair variables are
already available from the information obtained in round two.
For a pair at distance two, both endpoints know its common neighbors
and can choose the one with smallest identifier.  All messages and
enumerations in this schedule are finite; their size need not be
bounded.  The schedule applies to every finite graph, without knowing
its order or maximum degree and without testing planarity.


\paragraph{Simulating the packing incidence network.}
We apply the packing lemma to the residual vertex set $R$ and the
hypergraph $\mathcal H$.  The properties needed for this application are
\[
 \begin{gathered}
 1\leq |e|\leq2,\qquad \operatorname{diam}_G(e)\leq2
       \quad(e\in\mathcal H),\\
 |\{e\in\mathcal H:v\in e\}|\leq T,\qquad
 |N(v)\cap R|\leq B_0=Q(3Q+2)\quad(v\in V(G)).
 \end{gathered}
\]
The variables are $y_w$ for $w\in R$ and $z_e$ for
$e\in\mathcal H$.  A vertex row has the form
\begin{equation}\label{pack:mixed-row}
 \ell_v=\sum_{w\in R\cap N[v]}y_w
             +\sum_{e\in\mathcal H:v\in e}z_e\leq1.
\end{equation}
In addition there is one personal row $y_w\leq1/7$ for each $w\in R$.
All variables are nonnegative.  Add the common cap $M=1$.
This cap is redundant for $y_w$ because of its personal row, and for
$z_e$ because $e\ne\varnothing$: a unit-capacity row at any endpoint
of $e$ contains $z_e$, and all other terms in that row are nonnegative.
The vertex rows have at most $B_0+1+T$ variables and the personal rows
have one.  Thus Lemma~\ref{pack:theorem} applies with the fixed parameters
\begin{equation}\label{pack:parameters}
 K=B_0+1+T,\qquad C=\{1,1/7\},\qquad M=1.
\end{equation}
Let $L=|\mathcal P(K,C,1)|-1$ for these parameters.

Host the variable $y_w$, its personal row, and the vertex row at $w$
at the graph vertex $w$.  Host $z_e$ at the endpoint of $e$ with
smallest identifier.  Vertex row $v$ is hosted at $v$, whether or not
$v\in R$.  Every incidence of a $y$ variable with a vertex row has
host distance at most one, and the personal row is at the same host.
Every incidence of $z_e$ with a vertex row has host distance at most
two, by the diameter property of $e$.

The local preprocessing supplies these incidences and their endpoint
identifiers.  For a pair at distance two, choose a length-two route
deterministically, for example through its common neighbor of smallest
identifier; use its edge if the two endpoints are adjacent.  Constant
radius neighborhood information suffices to choose and record these
routes.  A message of the virtual incidence network is tagged with its
variable and row names and sent along the chosen route.  Every virtual
communication phase takes at most two graph rounds.  If many virtual
messages use the same physical edge, concatenate them into one finite
message.  This is permitted in LOCAL and creates no additional rounds.
Messages between nodes hosted at the same vertex require no communication.
Consequently the two phases of each palette level take at most four
graph rounds, for a total of at most $4L$ rounds for the packing scan.

At the end of the scan, each vertex row host knows its exact load
$\ell_v$, and each vertex knows its own $y_v$, with $y_v=0$ outside $R$.
It can decide locally whether
\[
 v\in S
 \quad\Longleftrightarrow\quad
 \ell_v\geq1-\tfrac12(1/7-y_v).
\]
After one exchange of these membership bits, each vertex can decide
whether it belongs to $U=R\setminus N[S]$ and hence to $O=S\cup U$.
Using at most ten rounds for the constant-radius preprocessing, and
allowing another ten rounds for route initialization and final
membership exchanges, the safe total bound is $4L+20$.  Its numerical
value is a single constant, since $Q$ and $T$ are fixed.  No vertex is
given $n$, the maximum degree, or a planar embedding.

All preprocessing tests enumerate finite vertex sets within fixed
radius neighborhoods, with the thresholds $Q$ and $T$ fixed.  Although
these neighborhoods may have unbounded size, they are finite, so the
tests are computable using the unbounded local computation allowed by
the model.  The rational packing arithmetic can be put on the common
integer grid described above.  If $Y_v=Dy_v$ and $L_v=D\ell_v$, the
selection test is the integer inequality
\[
 13D+7Y_v\leq14L_v.
\]
Thus the construction is an actual deterministic LOCAL program with
finite descriptions of its constants.  The same fixed program is
defined on every finite input graph; its correctness and approximation
guarantees use the planar preprocessing bounds.  It never needs to
test whether its input is planar.

\section{The slack available in rounding}
\label{sec:rounding}

The analysis compares the output with an arbitrary dominating set \(D\).
The algorithm never computes or receives this set. We first record the
slack furnished by rounding, and then show that almost all vertices of
\(D\) are suitable endpoints for planar contractions.

\subsection{Finite level sums}

All weights are rational. Choose a positive integer \(h\), divisible by
\(14\), such that \(hy_v\) is an integer for every vertex. This common
denominator is used only in the analysis. For a function \(F\) on the
positive grid points and \(a\in\{1/14,1/7\}\), write
\begin{equation}\label{eq:level-average}
 \avg{F(t)}{a}:=\frac1h\sum_{j=1}^{ah}F(j/h).
\end{equation}
This is a finite sum, so it is linear and preserves inequalities.
For a constant \(c\), \(\avg{c}{a}=ac\). For
\(p\in[0,\theta]\) with \(hp\) integral,
\begin{equation}\label{eq:layer-cake}
 \avg{\ind_{\{t\leq p\}}}{\theta}=p.
\end{equation}
Indeed, exactly \(hp\) of the grid points \(j/h\), \(1\leq j\leq h\theta\),
are at most \(p\). Thus this notation is a finite version of the usual
layer-cake argument; no limiting or integration theorem is needed.
If \(F(t)\geq0\), then
\(\avg{F(t)}{1/14}\leq\avg{F(t)}{\theta}\).

Put
\begin{equation}\label{eq:rounding-quantities}
 \begin{gathered}
 Y=\sum_{v\in V}y_v,\qquad Z=\sum_{e\in\HH}z_e,\qquad
 \delta=\sum_{s\in S}(\theta-y_s),\\
 B=N[S]\setminus S,\qquad Y_B=y(B),\\
 W_t=S\cup\{v\in B:y_v\geq t\},\qquad
 A_t=\{s\in S:y_s\geq t\},\\
 J=\avg{\Def(W_t)+\Def(A_t)}{\theta},\qquad
 J_0=\avg{\Def(W_t)}{\theta}.
 \end{gathered}
\end{equation}
All these quantities are nonnegative, and \(J_0\leq J\).
Every vertex with positive weight lies in exactly one of \(S,B,U\):
weights vanish off \(R\), and \(R\setminus N[S]=U\). By
\cref{lem:output-basic}, \(y(U)=\theta|U|\) and \(S\cap U=\varnothing\).
Also \(y(S)=\theta|S|-\delta\). Consequently
\begin{equation}\label{eq:mass-identity}
 7Y-|O|=7Y_B-7\delta.
\end{equation}

\begin{lemma}[Rounding with retained slack]\label{lem:rounding}
The algorithm's output satisfies
\begin{equation}\label{eq:rounding-slack}
 7Y-|O|\geq
 -\frac{13}{3}Z+\frac7{12}\delta+\frac12Y_B+\frac{13}{6}J.
\end{equation}
\end{lemma}
\begin{proof}
Consider the weighted neighbor incidence sum
\[
 M_S=\sum_{s\in S}\sum_{v\in N(s)}y_v.
\]
The selection condition in \eqref{eq:output} and the row definition give
\[
 \sum_{v\in N(s)}y_v
 \geq1-\theta+\frac{\theta-y_s}{2}
       -\sum_{\substack{e\in\HH\\s\in e}}z_e .
\]
Each hyperedge has at most two endpoints, so summing yields
\begin{equation}\label{eq:incidence-lower}
 M_S\geq\frac67|S|+\frac12\delta-2Z.
\end{equation}

For the upper bound, count the contribution of each graph edge to
\(\avg{e_G(W_t)+e_G(A_t)}{\theta}\).
An edge \(uv\) with both endpoints in \(S\) contributes
\(\theta+\min\{y_u,y_v\}\), which is at least \(y_u+y_v\), since
\(\max\{y_u,y_v\}\leq\theta\). This pays its two contributions to \(M_S\).
An edge from \(S\) to a vertex \(v\in B\) contributes exactly \(y_v\),
paying its sole contribution to \(M_S\). All other edge contributions
are nonnegative, and these two kinds account for every term in \(M_S\).
Using \eqref{eq:layer-cake} and the definition of the deficit, we get
\begin{align}
 M_S
 &\leq\avg{e_G(W_t)+e_G(A_t)}{\theta}\notag\\
 &=3\avg{|W_t|+|A_t|}{\theta}-J\notag\\
 &=3\bigl(\theta|S|+Y_B+y(S)\bigr)-J\notag\\
 &=\frac67|S|+3Y_B-3\delta-J.
 \label{eq:incidence-upper}
\end{align}
Comparing \eqref{eq:incidence-lower} and \eqref{eq:incidence-upper}
cancels the \(|S|\) term and proves
\[
 3Y_B\geq\frac72\delta+J-2Z.
\]
Substitute this in \eqref{eq:mass-identity} to obtain
\[
 7Y-|O|\geq-\frac{14}{3}Z+\frac76\delta+\frac73J.
\]
Take \(13/14\) of this inequality and add \(1/14\) of the exact
identity \eqref{eq:mass-identity}. The left side is unchanged. The
coefficients on the right become respectively
\(-13/3,\ 7/12,\ 1/2,\ 13/6\), proving the assertion.
\end{proof}

\subsection{Comparison with an integral dominating set}

Fix any dominating set \(D\). Separate the hyperedge mass into
\[
 Z_h=\sum_{\substack{e\in\HH\\e\cap D\ne\varnothing}}z_e,\qquad
 Z_m=\sum_{\substack{e\in\HH\\e\cap D=\varnothing}}z_e,
 \qquad Z=Z_h+Z_m .
\]
Here the subscripts mean \emph{hit} and \emph{missed}. Define the unused
capacity and the two overcount terms by
\begin{equation}\label{eq:comparison-slack}
 \begin{aligned}
 L_D&=\sum_{d\in D}(1-\ell_d),\\
 M_y&=\sum_{w\in R}y_w\bigl(|N[w]\cap D|-1\bigr),\\
 M_z&=\sum_{\substack{e\in\HH\\e\cap D\ne\varnothing}}
                  z_e\bigl(|e\cap D|-1\bigr),\\
 K_D&=L_D+M_y+M_z .
 \end{aligned}
\end{equation}
All terms are nonnegative: rows are feasible, \(D\) dominates every
vertex, and a hit hyperedge has at least one endpoint in \(D\).
Interchanging finite sums in the rows of \(D\) gives
\[
 \sum_{d\in D}\ell_d
 =\sum_{w\in R}y_w|N[w]\cap D|
   +\sum_{e\in\HH}z_e|e\cap D|
 =Y+M_y+Z_h+M_z .
\]
It follows that
\begin{equation}\label{eq:comparison-identity}
 |D|=Y+Z_h+K_D.
\end{equation}
This is an exact accounting identity, not an assertion that the packing
is optimal.

Define the common budget
\begin{equation}\label{eq:budget}
 \budget=K_D+Z_h+\delta+Y_B+J_0\geq0.
\end{equation}
Combining \cref{lem:rounding} with \eqref{eq:comparison-identity} gives
\begin{align}
 7|D|-|O|
 &\geq7K_D+\frac83Z_h+\frac7{12}\delta
          +\frac12Y_B+\frac{13}{6}J-\frac{13}{3}Z_m\notag\\
 &\geq\frac12\budget-\frac{13}{3}Z_m.
 \label{eq:budget-gap}
\end{align}
For the second inequality, use \(J\geq J_0\) and the fact that each
of \(7,8/3,7/12,1/2,13/6\) is at least \(1/2\).
We will prove that \(Z_m\leq\budget/100\).

\subsection{Exceptional vertices of the comparison set}

Let
\begin{equation}\label{eq:bad-good}
 D_b=\{d\in D:y_d<1/14\}\cup(D\cap U),\qquad
 D_g=D\setminus D_b,\qquad
 S_0=\{s\in S:y_s=0\}.
\end{equation}
The following bounds make these exceptional sets affordable.

\begin{lemma}\label{lem:bad-counts}
For every dominating set \(D\),
\begin{equation}\label{eq:bad-counts}
 |D_b|\leq35\budget,\qquad |S_0|\leq7\budget .
\end{equation}
Moreover, \(D_g\subseteq W_t\) and \(y_d>0\) for every \(d\in D_g\)
and every \(0<t\leq1/14\).
\end{lemma}
\begin{proof}
If \(s\in S\) has \(y_s<1/14\), its contribution
\(\theta-y_s\) to \(\delta\) is greater than \(1/14\).
If \(d\in D\setminus S\) has \(y_d<1/14\), the failure of the
selection condition gives
\[
 1-\ell_d>\frac{\theta-y_d}{2}>\frac1{28}.
\]
Thus
\begin{equation}\label{eq:low-weight-count}
 |\{d\in D:y_d<1/14\}|\leq14\delta+28L_D.
\end{equation}

For \(d\in D\cap U\), there is no member of \(S\) in \(N[d]\).
There are at most six members of \(U\) in \(N[d]\): seven such
vertices would, by \(y_u=1/7\), force \(\ell_d\geq1\), making
\(d\in S\), a contradiction. Every positive-weight vertex in
\(N[d]\) therefore belongs to \(U\) or \(B\), and
\[
 \ell_d
 \leq\frac67+\sum_{w\in B\cap N[d]}y_w
              +\sum_{\substack{e\in\HH\\d\in e}}z_e.
\]
Rearrange and sum over \(d\in D\cap U\).
The unused row capacities sum to at most \(L_D\).
The boundary terms sum to at most \(Y_B+M_y\): for each \(w\in B\),
its multiplicity is at most \(|N[w]\cap D|\), and a positive
weight implies \(w\in R\). The hyperedge terms sum to at most
\(Z_h+M_z\), by the definition of \(M_z\).
Consequently
\begin{equation}\label{eq:du-count}
 |D\cap U|\leq7(L_D+Y_B+M_y+Z_h+M_z)
             =7(K_D+Y_B+Z_h).
\end{equation}
Combining \eqref{eq:low-weight-count} and \eqref{eq:du-count}, and
using \(L_D\leq K_D\), gives
\[
 |D_b|\leq14\delta+35K_D+7Y_B+7Z_h\leq35\budget.
\]
Each member of \(S_0\) contributes \(1/7\) to \(\delta\), so
\(|S_0|\leq7\delta\leq7\budget\).

Finally, a member \(d\in D_g\) has \(y_d\geq1/14>0\), hence lies
in \(R\). Since it is not in \(U=R\setminus N[S]\), it lies in
\(N[S]\). It therefore belongs either to \(S\) or to
\(\{v\in B:y_v\geq t\}\) for every \(0<t\leq1/14\).
This proves \(D_g\subseteq W_t\).
\end{proof}

\section{A planar count for many pairs}
\label{sec:faces}

\begin{lemma}[Nontriangular darts, including small components]
\label{geom:nontriangular-darts}
Let $J$ be a finite simple planar graph. Draw its connected components
separately on spheres, and use their facial permutations
$\varphi=\sigma\circ\alpha$ as above. Let $b$ be the number of darts on
facial boundaries whose length is not three. Then, with
$\Delta_J=3|V(J)|-|E(J)|$,
\begin{equation}\label{geom:bad-dart-bound}
 b\leq4\Delta_J.
\end{equation}
\end{lemma}
\begin{proof}
For a connected component with $n\geq3$ vertices, all facial boundaries
have length at least three. A nontriangular one has length $\ell\geq4$,
so $\ell\leq4(\ell-3)$. Summing, and including zero contributions from
triangles, gives
\[
 b_{\rm comp}
 \leq4\sum_{\rm faces}(\ell-3)
 =4(2e-3f)=4(3n-e-6)\leq4(3n-e).
\]
The equality uses Euler's formula. A component on two vertices is $K_2$.
Its two darts form one boundary of length two, so
$b_{\rm comp}=2\leq4(6-1)$. An isolated vertex has no darts and satisfies
$b_{\rm comp}=0\leq4\cdot3$. Sum over components; both $b$ and $3n-e$
are additive. The empty graph has both sides zero.
\end{proof}

\begin{figure}[tbp]
\centering
\begin{tikzpicture}[scale=1.1]
 \begin{scope}
  \coordinate (a) at (-1.2,-.6);
  \coordinate (b) at (1.2,-.6);
  \coordinate (c) at (0,1.5);
  \coordinate (x) at (0,.15);
  \fill[orange!15] (a)--(x)--(b)--cycle;
  \draw[edge] (a)--(c)--(b) (a)--(x)--(b) (c)--(x);
  \draw[newedge,-{Stealth[length=2.4mm]},
    shorten >=3.5mm,shorten <=3.5mm] (b)--(a);
  \node[vertex] at (a) {\(a\)};
  \node[vertex] at (b) {\(b\)};
  \node[vertex] at (c) {\(c\)};
  \node[witness] at (x) {\(x\)};
  \node[font=\small,align=center] at (0,-1.9)
   {A triangular corner identifies \(x\)\\from the oriented side \((b,a)\)};
 \end{scope}
 \begin{scope}[xshift=6cm]
  \coordinate (a) at (-1.2,-.6);
  \coordinate (b) at (1.2,-.6);
  \coordinate (c) at (0,1.5);
  \coordinate (d) at (0,-1.35);
  \coordinate (x) at (0,.15);
  \fill[orange!15] (x)--(a)--(d)--(b)--cycle;
  \draw[edge] (a)--(c)--(b) (a)--(x)--(b) (c)--(x)
    (a)--(d)--(b);
  \draw[thick,dashed,pairpurple] (a)--(b);
  \node[vertex] at (a) {\(a\)};
  \node[vertex] at (b) {\(b\)};
  \node[vertex] at (c) {\(c\)};
  \node[vertex] at (d) {\(d\)};
  \node[witness] at (x) {\(x\)};
  \node[font=\small,align=center] at (0,-1.9)
   {A nontriangular corner is charged\\to the planar deficit};
 \end{scope}
\end{tikzpicture}
\caption{The two types of degree-three witnesses in the face count.
Blue vertices belong to \(W\); the orange vertex \(x\) lies outside \(W\).
On the left, the face with boundary \(a,x,b,a\) assigns the pair
incidence to the dart \((b,a)\), whose next facial dart ends at \(x\).
On the right, the shaded face has boundary \(x,a,d,b,x\).
The dashed segment \(ab\) denotes a pair in \(\mathcal P\), not an
edge of the planar graph. This distinction is allowed in the lemma.}
\label{fig:corners}
\end{figure}

\begin{lemma}[Triangular corners of degree-three vertices]
\label{geom:triangle-injection}
Fix planar rotations of $J$. Let $\mathcal P$ be any graph of pairs on
$V(J)$ with $m$ edges. Let $C_3$ consist of vertices of degree three in
$J$ for which every incident corner belongs to a triangular face. Then
\[
 |\{(x,\{a,b\}):x\in C_3,\ \{a,b\}\in E(\mathcal P),
                         \ a,b\in N_J(x)\}|\leq2m.
\]
\end{lemma}
\begin{proof}
Represent a corner at $x$ by an incoming dart $q=(a,x)$. If $(x,b)$
follows $(x,a)$ in the cyclic order at $x$, then $\varphi q=(x,b)$.
The corner is triangular precisely when $q$ lies in a face cycle of
length three.

Since $x$ has three neighbors, their cyclic successor is a three-cycle.
For any distinct neighbors $a,b$, exactly one follows the other.
Orient their names so that $a$ is followed by $b$, and take $q=(a,x)$.
The triangular corner gives
\[
 q=(a,x),\qquad \varphi q=(x,b),\qquad
 p:=\varphi^2q=(b,a),\qquad \varphi p=q.
\]
In particular, the opposite side $ab$ is an edge of $J$.
Map the incidence $(x,\{a,b\})$ to this oriented side $p=(b,a)$.
The dart $p$ recovers its underlying unordered pair and recovers $x$ as
the final endpoint of $\varphi p$. Thus the map is injective. Each of
the $m$ edges of $\mathcal P$ has only two orientations. A pair that is
not an edge of $J$ cannot occur as an image. The asserted bound follows.
\end{proof}

\begin{lemma}[Ambient pair incidence]\label{geom:ambient-pairs}
Let $W,C$ be disjoint vertex sets in a finite simple planar graph $G$.
Suppose every $x\in C$ has at least three neighbors in $W$.
Let $\mathcal P$ be a graph of pairs on $W$ with $m$ edges and maximum
degree at most $T$; its edges need not belong to $G$. Then
\begin{equation}\label{geom:ambient-strong}
 \sum_{\{a,b\}\in E(\mathcal P)}|N_G(a)\cap N_G(b)\cap C|
 \leq2m+6T\Def(W).
\end{equation}
In particular, the inequality holds with $12T$ in place of $6T$.
No connectivity or minimum-order assumption is required.
\end{lemma}
\begin{proof}
Let $J$ have vertex set $W\cup C$, all edges of $G[W]$, and all edges
between $W$ and $C$. This deletes any edges inside $C$, so $J$ is
planar, $C$ is independent, and
$d_x:=\deg_J(x)=|N_G(x)\cap W|\geq3$ for $x\in C$.
The incidence count is unchanged; by interchanging finite sums it is
\[
 I=\sum_{x\in C}c_J(x),\qquad
 c_J(x)=|\{\{a,b\}\in E(\mathcal P):a,b\in N_J(x)\}|.
\]
Set
\[
 E_C=\sum_{x\in C}(d_x-3)\geq0,\qquad
 \Delta_J=3|V(J)|-|E(J)|\geq0.
\]
Every edge of $J$ is either inside $W$ or has exactly one endpoint in
$C$. Hence
\[
 |E(J)|=e_G(W)+\sum_{x\in C}d_x,
 \qquad
 \Def(W)=\Delta_J+E_C.                         \tag{*}
\]

Fix componentwise planar rotations, and let $b$ be the number of
nontriangular darts. Partition $C$ into: degree-three vertices whose
corners are all triangular; degree-three vertices having a nontriangular
corner; and vertices of degree at least four. The first class contributes
at most $2m$ by Lemma~\ref{geom:triangle-injection}.

For each vertex $x$ of the second class, choose an incoming dart
$q_x=(a_x,x)$ in a nontriangular corner. These darts are distinct,
because their final endpoints are distinct. Thus there are at most $b$
vertices in this class. Each contributes at most $3T/2$, by
Lemma~\ref{geom:pair-degree} applied to its three neighbors.

For $d_x\geq4$, the same bound gives
\[
 c_J(x)\leq Td_x/2\leq2T(d_x-3).
\]
The last inequality follows from $d_x\geq4$ and $T\geq0$. Summing it
bounds the third class by $2TE_C$. Combining the three bounds and
Lemma~\ref{geom:nontriangular-darts} yields
\[
 \begin{split}
 I&\leq2m+\tfrac32Tb+2TE_C\\
  &\leq2m+6T\Delta_J+2TE_C\\
  &\leq2m+6T(\Delta_J+E_C)=2m+6T\Def(W),
 \end{split}
\]
where the last equality is (*). The $12T$ version follows from
$\Def(W)\geq0$.
\end{proof}

\section{Charging hyperedges missed by a dominating set}
\label{sec:charging}

Fix an arbitrary dominating set $D$, and a domination assignment
\[
 f:V(G)\to D,\qquad f(w)\in N[w],\qquad f(d)=d\quad(d\in D).
\]
This choice is only for analysis; the algorithm does not know $D$ or
$f$. Let $m_1,m_2$ count the singleton and pair edges of $\mathcal H$
missed by $D$. Let $\mathcal P$ be the graph of the $m_2$ missed pairs;
its maximum degree is at most $T$.

For $0<t\leq1/14$, recall
\[
 W_t=S\cup\{v\in N[S]\setminus S:y_v\geq t\},\qquad
 S_0=\{s\in S:y_s=0\},
\]
\[
 D_b=\{d\in D:y_d<1/14\}\cup(D\cap U),\qquad D_g=D\setminus D_b.
\]
The geometric arguments use
\begin{equation}\label{geom:good-in-level}
 D_g\subseteq W_t,\qquad y_d>0\quad(d\in D_g),
\end{equation}
and
\begin{equation}\label{geom:common-zero}
 w\in X(e),\ e\in\mathcal H\quad\Longrightarrow\quad y_w=0.
\end{equation}
Indeed, a member of $D_g$ has weight at least $1/14$, hence belongs to
$R$, and does not belong to $U=R\setminus N[S]$. It therefore lies in
$N[S]$ and in $W_t$. The second fact holds because
$X(e)\subseteq X=V(G)\setminus R$ and the weights are zero outside $R$.

Every retained hyperedge meets $S$. If all endpoint rows of a hyperedge
were slack, its $z_e$ coordinate could be increased, contrary to
maximality. If the redundant cap $z_e\leq1$ is attained, its unit rows
are already tight. A tight row belongs to $S$, because $y_v\leq1/7$.
In particular, every retained singleton center is in $S$.

\begin{lemma}[Singleton charge]\label{geom:singleton-charge}
For every $0<t\leq1/14$,
\begin{equation}\label{geom:singleton-charge-eq}
 (T-9)m_1\leq2|S_0|+2|D_b|+\Def(W_t).
\end{equation}
\end{lemma}
\begin{proof}
Let $A$ be the missed singleton centers. Then $|A|=m_1$,
$A\cap D=\varnothing$, and $A\subseteq S\subseteq W_t$.
Each $\{a\}$ is $T$-hard. The set $f(N(a))\subseteq D$ dominates
$N(a)$, since each $w\in N(a)$ has its assigned dominator in this set.
It avoids $a$, so hardness gives $|f(N(a))|\geq T+1$.
Consequently
\[
 P=\{(a,d):a\in A,\ d\in f(N(a))\}
 \quad\text{satisfies}\quad |P|\geq(T+1)m_1.
\]
For each distinct pair $(a,d)\in P$, choose one witness
$w=w_{a,d}\in N(a)$ with $f(w)=d$. The map
\begin{equation}\label{geom:witness-injection}
 (a,d)\longmapsto(a,w_{a,d})
\end{equation}
is injective: the image recovers $a$ and $w$, hence $d=f(w)$.

Partition $P$ in the following order:
\begin{align*}
 P_1&=\{(a,d)\in P:w_{a,d}\in A\},\\
 P_2&=\{(a,d)\in P:w_{a,d}\notin A,\ w_{a,d}\in S_0\},\\
 P_3&=\{(a,d)\in P:w_{a,d}\notin A\cup S_0,\ d\in D_b\},\\
 P_4&=\{(a,d)\in P:w_{a,d}\notin A\cup S_0,\ d\in D_g\}.
\end{align*}
These classes are disjoint and exhaustive. Every chosen witness has
weight zero, since it belongs to $N(a)=X(\{a\})$ and
\eqref{geom:common-zero} applies. A witness outside $A\cup S_0$ is
therefore outside $S$ and, as $t>0$, outside $W_t$.

For $P_1$, \eqref{geom:witness-injection} maps into the oriented edges
of $G[A]$, so
\[
 |P_1|\leq2e_G(A)\leq6m_1.
\]
For $P_2$, it maps into edges between the disjoint sets $A$ and
$S_0\setminus A$, with one orientation per edge. Thus
\[
 |P_2|\leq e_G(A,S_0\setminus A)
 \leq2(|A|+|S_0\setminus A|)\leq2(m_1+|S_0|).
\]
For $P_3$, Lemma~\ref{geom:assigned-incidences}, with endpoint sets
$A,D_b$, applies because each witness avoids $A$. It gives
\[
 |P_3|\leq2(m_1+|D_b|).
\]
A witness already in $D_b$ is allowed: it is fixed as its own assigned
center and is not contracted to another vertex.

For $P_4$, both endpoints are in $W_t$, while each witness is outside
$W_t$. The edge $ad$ is absent from $G$: otherwise $d\in N(a)$ would
have weight zero by \eqref{geom:common-zero}, contradicting
$d\in D_g$ and \eqref{geom:good-in-level}. The new-edge assertion of
Lemma~\ref{geom:assigned-incidences} therefore gives
\[
 |P_4|\leq\Def(W_t).
\]
All contractions for this last class can be performed together while
fixing $W_t$ and retaining its induced edges. They add the distinct
previously absent edges indexed by $P_4$. The contraction used for
$P_3$ need not be the same contraction.

Adding the four estimates proves
\[
 (T+1)m_1\leq|P|
 \leq10m_1+2|S_0|+2|D_b|+\Def(W_t).
\]
Subtracting $10m_1$ gives \eqref{geom:singleton-charge-eq}.
\end{proof}

\begin{lemma}[Pair charge]\label{geom:pair-charge}
For every $0<t\leq1/14$,
\begin{equation}\label{geom:pair-charge-eq}
 (Q-86T-2)m_2\leq3T|D_b|+T|S_0|+12T\Def(W_t).
\end{equation}
\end{lemma}
\begin{proof}
Let $A$ be the endpoints of the $m_2$ missed pairs. Then
$|A|\leq2m_2$ and $A\cap D=\varnothing$.
Every missed pair $e$ is $Q$-hard. Its common neighborhood $X(e)$
avoids $e$, because the graph has no loops, and dominates itself using
closed neighborhoods. Hence $|X(e)|\geq Q+1$.
The incidence set
\[
 \mathcal I=\{(e,w):e\in E(\mathcal P),\ w\in X(e)\}
\]
satisfies
\begin{equation}\label{geom:pair-incidence-lower}
 |\mathcal I|=\sum_{e\in E(\mathcal P)}|X(e)|
 \geq(Q+1)m_2\geq Qm_2.
\end{equation}
With $c(w)$ as in Lemma~\ref{geom:pair-degree},
$|\mathcal I|=\sum_w c(w)$. Every witness with $c(w)>0$ has weight
zero by \eqref{geom:common-zero}.

Partition these incidences by their witness:
\begin{enumerate}[label=(\roman*)]
\item $w\in A$;
\item $w\notin A$ and $w\in S_0$;
\item $w\notin A\cup S_0$ and $f(w)\in D_b$;
\item $w\notin A\cup S_0$ and $f(w)\in D_g$.
\end{enumerate}
These classes are disjoint and exhaustive. Witnesses in the last two
classes lie outside $W_t$: they have weight zero, and membership in
$S$ would put them in $S_0$.

By \eqref{geom:pairs-inside}, the first class has size at most
\begin{equation}\label{geom:pair-cat-one}
 3T|A|\leq6Tm_2.
\end{equation}
By \eqref{geom:pairs-bipartite} for $A,S_0\setminus A$, the second
class has size at most
\begin{equation}\label{geom:pair-cat-two}
 T(|A|+|S_0\setminus A|)\leq2Tm_2+T|S_0|.
\end{equation}

To treat multiplicities in the third class, let $C_b$ be its set of
distinct witnesses. Contract the vertices of $C_b\setminus D$ to their
assigned centers as in Lemma~\ref{geom:star-contraction}, obtaining a
simple planar graph $F$. Witnesses already in $D$ are left fixed as
their own centers. Every vertex of $A$ and $D_b$ is fixed, since $C_b$
avoids $A$ and no vertex of $D$ is a leaf.

For an incidence $(e,w)$ in this class, write $e=\{a,b\}$ and
$d=f(w)$. The three vertices $a,b,d$ are distinct, because $a,b$ are
distinct members of $A$ and $d\in D$. The edges $aw,bw$ become $ad,bd$
in $F$. Map the incidence to $(e,d)$. For any fixed image
$(\{a,b\},d)$, its preimages are distinct members of
\[
 N_G(a)\cap N_G(b)\cap N_G[d].
\]
Lemma~\ref{geom:outside-pair} bounds their number by three, including
the possible preimage $w=d$.

Every image $(e,d)$ has both endpoints of $e$ in $N_F(d)\cap A$.
The pair graph is still the original $\mathcal P$, with maximum
degree at most $T$. Applying \eqref{geom:pairs-bipartite} in the planar
graph $F$ to the disjoint sets $A,D_b$ bounds the number of possible
images by $T(|A|+|D_b|)$. Multiplying by the fiber bound three gives
\begin{equation}\label{geom:pair-cat-three}
 |\mathcal I_{\rm iii}|
 \leq3T(|A|+|D_b|)\leq6Tm_2+3T|D_b|.
\end{equation}
Any extra edges created by contraction only enlarge the set of possible
images; that enlarged set is still bounded by the planar bipartite
edge count.

For the fourth class, put $W'=W_t\cup A$ and let $C_g$ be its set of
distinct witnesses. Each $w\in C_g$ lies outside $W'$, while
$d=f(w)\in D_g\subseteq W_t$. Choose any missed pair $\{a,b\}$
witnessed by $w$. The three vertices $a,b,d$ are distinct and in $W'$.
The edges $wa,wb$ exist. Also $wd$ exists: domination gives either
$w=d$ or adjacency, and equality is impossible because $w\notin W'$
and $d\in W'$. Hence $w$ has at least three neighbors in $W'$.
All pair endpoints are in $W'$. The $12T$ form of
Lemma~\ref{geom:ambient-pairs} gives
\begin{equation}\label{geom:pair-cat-four-pre}
 |\mathcal I_{\rm iv}|\leq2m_2+12T\Def(W').
\end{equation}
This counts precisely the incidences at witnesses in $C_g$, since
their category depends only on $w$, not on which pair they witness.

Adding the vertices of $A$ retains every old induced edge on $W_t$.
Therefore
\[
 \begin{split}
 \Def(W')&=3|W'|-e_G(W')\\
 &\leq3(|W_t|+|A|)-e_G(W_t)
 =\Def(W_t)+3|A|\leq\Def(W_t)+6m_2.
 \end{split}
\]
Thus
\begin{equation}\label{geom:pair-cat-four}
 |\mathcal I_{\rm iv}|\leq2m_2+12T\Def(W_t)+72Tm_2.
\end{equation}
Combining \eqref{geom:pair-incidence-lower} with the four bounds gives
\[
 Qm_2\leq(86T+2)m_2+3T|D_b|+T|S_0|+12T\Def(W_t),
\]
where $86=6+2+6+72$. Rearrangement proves the claim. For $m_2=0$,
all the incidence sets are empty and the same inequalities remain valid.
\end{proof}

\section{The exact endpoint}
\label{sec:endpoint}

We now combine the pointwise planar charges with the budget from
\cref{sec:rounding}. Recall that \(m_1,m_2\) count the hyperedges
missed by the arbitrary dominating set \(D\).

\begin{lemma}\label{lem:missed-budget}
For the fixed constants in \eqref{eq:constants},
\[
 Z_m\leq\frac{\budget}{100}.
\]
\end{lemma}
\begin{proof}
Apply the finite average \eqref{eq:level-average} on
\(0<t\leq1/14\) to \eqref{geom:singleton-charge-eq}, and multiply
by \(14\). Since the deficit is nonnegative,
\(\avg{\Def(W_t)}{1/14}\leq J_0\), and hence
\[
 (T-9)m_1\leq2|S_0|+2|D_b|+14J_0
 \leq(14+70+14)\budget=98\budget.
\]
Here we used \cref{lem:bad-counts} and \(J_0\leq\budget\).
The same operation on \eqref{geom:pair-charge-eq} gives
\[
 \begin{split}
 (Q-86T-2)m_2
 &\leq3T|D_b|+T|S_0|+168TJ_0\\
 &\leq(105+7+168)T\budget=280T\budget.
 \end{split}
\]
Both denominators below are strictly positive. More precisely,
\[
 T-9=999991>98000,\qquad
 Q-86T-2=999913999998>280000T=280000000000.
\]
Therefore
\[
 \frac{98}{T-9}<\frac1{1000},\qquad
 \frac{280T}{Q-86T-2}<\frac1{1000}.
\]
Every \(z_e\leq1\), so
\[
 Z_m\leq m_1+m_2
 \leq\left(\frac{98}{T-9}+
           \frac{280T}{Q-86T-2}\right)\budget
 \leq\frac{\budget}{100}.
\]
The last weak inequality follows from
\(2/1000<1/100\) and \(\budget\geq0\); it also covers
\(\budget=0\) without a separate strict-inequality argument.
\end{proof}

\begin{proof}[Proof of \cref{thm:main}]
The construction of \cref{sec:algorithm} is an effective deterministic
constant-round algorithm, by \cref{pack:theorem} and the implementation
in \cref{sec:implementation}. Its output \(O\) dominates \(G\) by
\cref{lem:output-basic}. For every dominating set \(D\),
\eqref{eq:budget-gap} and \cref{lem:missed-budget} yield
\[
 7|D|-|O|
 \geq\frac12\budget-\frac{13}{3}Z_m
 \geq\left(\frac12-\frac{13}{300}\right)\budget
 =\frac{137}{300}\budget\geq0.
\]
A minimum dominating set exists because the graph is finite and the
whole vertex set is dominating. Choosing such a \(D\) proves
\(|O|\leq7\gamma(G)\). All estimates used relaxed planar edge bounds
valid for every order and every number of components, so no exceptional
small or disconnected cases need to be removed. On the empty graph
the algorithm outputs the empty set and the same inequality holds.
\end{proof}

The number of rounds is one fixed integer, at most \(4L+20\) for the
explicit finite palette specified in \cref{sec:implementation}.
Neither a limit over algorithms nor an input-dependent choice of an
error tolerance is involved. The packing is only coordinatewise
maximal, and its exact computation uses bounded row size rather than
a maximum-degree bound on the input graph. The final comparison
uses an integral dominating set through its domination assignment;
the theorem does not assert the stronger inequality
\(|O|\leq7\gamma_f(G)\).

Combined with the lower bound of Hilke, Lenzen, and Suomela~\cite{HLS14},
the theorem proves that seven is both the infimum and an attained
approximation factor for deterministic constant-round planar domination.

\begin{thebibliography}{99}
\small\raggedright

\bibitem{Linial92}
Nathan Linial.
Locality in distributed graph algorithms.
\emph{SIAM Journal on Computing}, 21(1):193--201, 1992.
\href{https://doi.org/10.1137/0221015}{doi:10.1137/0221015}.

\bibitem{Suomela13}
Jukka Suomela.
Survey of local algorithms.
\emph{ACM Computing Surveys}, 45(2), Article 24, 1--40, 2013.
\href{https://doi.org/10.1145/2431211.2431223}{doi:10.1145/2431211.2431223}.

\bibitem{LPW13}
Christoph Lenzen, Yvonne-Anne Pignolet, and Roger Wattenhofer.
Distributed minimum dominating set approximations in restricted families
of graphs.
\emph{Distributed Computing}, 26(2):119--137, 2013.
\href{https://doi.org/10.1007/s00446-013-0186-z}{doi:10.1007/s00446-013-0186-z}.

\bibitem{Waw14}
Wojciech Wawrzyniak.
A strengthened analysis of a local algorithm for the minimum dominating
set problem in planar graphs.
\emph{Information Processing Letters}, 114(3):94--98, 2014.
\href{https://doi.org/10.1016/j.ipl.2013.11.008}{doi:10.1016/j.ipl.2013.11.008}.

\bibitem{HLS14}
Miikka Hilke, Christoph Lenzen, and Jukka Suomela.
Brief announcement: Local approximability of minimum dominating set on
planar graphs.
In \emph{Proceedings of PODC 2014}, pages 344--346, 2014.
\href{https://doi.org/10.1145/2611462.2611504}{doi:10.1145/2611462.2611504}.
Full version: \href{https://arxiv.org/abs/1402.2549}{arXiv:1402.2549}.

\bibitem{HSV22}
Ozan Heydt, Sebastian Siebertz, and Alexandre Vigny.
Local planar domination revisited.
In \emph{SIROCCO 2022}, Lecture Notes in Computer Science 13298,
pages 154--173. Springer, 2022.
\href{https://doi.org/10.1007/978-3-031-09993-9_9}{doi:10.1007/978-3-031-09993-9\_9}.

\bibitem{HKSS25}
Ozan Heydt, Simeon Kublenz, Patrice Ossona de Mendez, Sebastian Siebertz,
and Alexandre Vigny.
Distributed domination on sparse graph classes.
\emph{European Journal of Combinatorics}, 123:103773, 2025.
\href{https://doi.org/10.1016/j.ejc.2023.103773}{doi:10.1016/j.ejc.2023.103773}.
\href{https://arxiv.org/abs/2207.02669}{arXiv:2207.02669}.

\bibitem{Waw26}
Wojciech Wawrzyniak.
Near-optimal distributed domination in planar graphs.
Preprint, 2026.
\href{https://arxiv.org/abs/2609.32736}{arXiv:2609.32736}.

\bibitem{AS26}
Saeed Amiri and Sebastian Siebertz.
Almost optimal constant-round approximation of dominating set in graph
classes with excluded minors.
Preprint, 2026.
\href{https://arxiv.org/abs/2610.09618}{arXiv:2610.09618}.

\bibitem{Picavet26}
Timoth{\'e} Picavet.
Almost tight distributed approximation for domination in sparse graph classes.
Preprint, 2026.
\href{https://arxiv.org/abs/2610.31203}{arXiv:2610.31203}.

\bibitem{BU17}
Nikhil Bansal and Seeun William Umboh.
Tight approximation bounds for dominating set on graphs of bounded arboricity.
\emph{Information Processing Letters}, 122:21--24, 2017.
\href{https://doi.org/10.1016/j.ipl.2017.01.011}{doi:10.1016/j.ipl.2017.01.011}.

\bibitem{Dvorak19}
Zden{\v e}k Dvo{\v r}{\'a}k.
On distance $r$-dominating and $2r$-independent sets in sparse graphs.
\emph{Journal of Graph Theory}, 91(2):162--173, 2019.
\href{https://doi.org/10.1002/jgt.22426}{doi:10.1002/jgt.22426}.
\href{https://arxiv.org/abs/1710.10010}{arXiv:1710.10010}.

\bibitem{DGI24}
Michal Dory, Mohsen Ghaffari, and Saeed Ilchi.
Near-optimal distributed dominating set in bounded arboricity graphs.
\emph{Distributed Computing}, 37(4):387--398, 2024.
\href{https://doi.org/10.1007/s00446-023-00447-z}{doi:10.1007/s00446-023-00447-z}.
\href{https://arxiv.org/abs/2206.05174}{arXiv:2206.05174}.

\bibitem{KMW06}
Fabian Kuhn, Thomas Moscibroda, and Roger Wattenhofer.
The price of being near-sighted.
In \emph{Proceedings of SODA 2006}, pages 980--989, 2006.
\href{https://doi.org/10.1145/1109557.1109666}{doi:10.1145/1109557.1109666}.

\bibitem{KMW16}
Fabian Kuhn, Thomas Moscibroda, and Roger Wattenhofer.
Local computation: Lower and upper bounds.
\emph{Journal of the ACM}, 63(2), Article 17, 1--44, 2016.
\href{https://doi.org/10.1145/2742012}{doi:10.1145/2742012}.
\href{https://arxiv.org/abs/1011.5470}{arXiv:1011.5470}.

\bibitem{AS10}
Matti {\AA}strand and Jukka Suomela.
Fast distributed approximation algorithms for vertex cover and set cover
in anonymous networks.
In \emph{Proceedings of SPAA 2010}, pages 294--302, 2010.
\href{https://doi.org/10.1145/1810479.1810533}{doi:10.1145/1810479.1810533}.

\bibitem{GHS17}
Mika G{\"o}{\"o}s, Juho Hirvonen, and Jukka Suomela.
Linear-in-$\Delta$ lower bounds in the LOCAL model.
\emph{Distributed Computing}, 30:325--338, 2017.
\href{https://doi.org/10.1007/s00446-015-0245-8}{doi:10.1007/s00446-015-0245-8}.
\href{https://arxiv.org/abs/1304.1007}{arXiv:1304.1007}.

\bibitem{Dahal24}
Sameep Dahal and Jukka Suomela.
Distributed half-integral matching and beyond.
\emph{Theoretical Computer Science}, 982:114278, 2024.
\href{https://doi.org/10.1016/j.tcs.2023.114278}{doi:10.1016/j.tcs.2023.114278}.
\href{https://arxiv.org/abs/2303.05250}{arXiv:2303.05250}.

\bibitem{FHK11}
Patrik Flor{\'e}en, Marja Hassinen, Joel Kaasinen, Petteri Kaski,
Topi Musto, and Jukka Suomela.
Local approximability of max-min and min-max linear programs.
\emph{Theory of Computing Systems}, 49(4):672--697, 2011.
\href{https://doi.org/10.1007/s00224-010-9303-6}{doi:10.1007/s00224-010-9303-6}.

\bibitem{ER60}
Paul Erd{\H o}s and Richard Rado.
Intersection theorems for systems of sets.
\emph{Journal of the London Mathematical Society}, s1-35(1):85--90, 1960.
\href{https://doi.org/10.1112/jlms/s1-35.1.85}{doi:10.1112/jlms/s1-35.1.85}.

\bibitem{Diestel25}
Reinhard Diestel.
\emph{Graph Theory}, sixth edition.
Graduate Texts in Mathematics 173. Springer, 2025.
\href{https://doi.org/10.1007/978-3-662-70107-2}{doi:10.1007/978-3-662-70107-2}.

\end{thebibliography}

\end{document}
